%=============================================================================%
%  福州大学 专业学位研究生论文开题报告 —— LaTeX 空白模板
%  版式照福州大学开题报告 Word 模板（FZU Body / FZU H1~H3 样式）
%
%  需要填写的只有三处：
%    1) contents/frontinfo_fzu.tex   封面信息（题目/姓名/学号/导师/专业/单位）
%    2) contents/mainmatter_fzu.tex  正文五章（骨架已给，注释里是每节的填写要点）
%    3) bib/ref.bib                  参考文献条目（与主模板共用文献池）
%
%  编译：xelatex proposal_fzu.tex && bibtex proposal_fzu && xelatex proposal_fzu.tex && xelatex proposal_fzu.tex
%=============================================================================%
\documentclass{style/ucasproposal}
\usepackage[super, list, xhf]{style/artratex}% xhf：关闭模板自带的页眉页脚，改用文末自定义的 fzu 样式
\usepackage{style/artracom}
\usepackage{float}% 支持 [H] 强制图片位置
\usepackage{multirow}% 封面表一的合并单元格
\usepackage{tabularx}% 封面表一按版心宽度自适应
%---------------------------------------------------------------------------%
% 版面与字体：按福大开题报告模板（表一 .docx 内嵌的 FZU Body / FZU H1~H3 样式）
%   页面   A4，上下 22mm、左右 20mm（pgMar 1247 / 1134 twips）
%   正文   宋体 + Times New Roman，小四 12pt，固定行距 20 磅，首行缩进 2 字符
%   标题1  黑体四号 14pt 加粗，段前段后各 12 磅，编号后留 1.5 字间隔
%   标题2  黑体小四 12pt，编号后留间隔
%   标题3  黑体小四 12pt，左对齐，段前 12 磅、段后 6 磅
%   图号   按章编号（图 1-1）；公式连续编号（1）
%   页脚   宋体小五 9pt，居中「第 X 页 共 Y 页」，无页眉
%---------------------------------------------------------------------------%
\usepackage{geometry}
\usepackage{fancyhdr}
\geometry{paper=a4paper,left=20mm,right=20mm,top=22mm,bottom=22mm,
          headheight=15mm,headsep=7mm,footskip=17.5mm}
% 字体：Windows 用宋体/黑体/仿宋；macOS 自动换宋体-简/楷体；都没有时用 Fandol 兜底，
%       保证在任何一台装了 TeX 发行版的机器上都不会因缺字体而报错。
\IfFontExistsTF{Times New Roman}{%
  \setmainfont{Times New Roman}%
}{%
  \setmainfont{TeX Gyre Termes}%
}
\IfFontExistsTF{SimSun}{% Windows
  \setCJKmainfont{SimSun}[AutoFakeBold=2.5,ItalicFont=KaiTi]
  \setCJKsansfont{SimHei}[AutoFakeBold=2.5]
  \setCJKmonofont{FangSong}
}{\IfFontExistsTF{Songti SC}{% macOS
  \setCJKmainfont{Songti SC}[BoldFont=* Bold,ItalicFont=Kaiti SC]
  \setCJKsansfont{Heiti SC}[BoldFont=* Medium]
  \setCJKmonofont{STFangsong}
}{% Linux / Overleaf
  \IfFileExists{FandolSong-Regular.otf}{% Fandol 随 TeX Live 发行，无需装字体
    \setCJKmainfont{FandolSong-Regular}[Extension=.otf,BoldFont=FandolSong-Bold,ItalicFont=FandolKai-Regular]
    \setCJKsansfont{FandolHei-Regular}[Extension=.otf,BoldFont=FandolHei-Bold]
    \setCJKmonofont{FandolFang-Regular}[Extension=.otf]
  }{% 最后兜底：Linux 上常见的 Noto CJK
    \setCJKmainfont{Noto Serif CJK SC}[BoldFont=* Bold]
    \setCJKsansfont{Noto Sans CJK SC}[BoldFont=* Bold]
    \setCJKmonofont{Noto Serif CJK SC}
  }
}}
\setlength{\parindent}{24pt}% 首行缩进 2 字符（480 twips）
\setlength{\parskip}{6pt plus 0pt minus 0pt}% 段后 6 磅（after=120 twips）
\ctexset{
  section = {
    format = \raggedright\linespread{1.0}\fontsize{14pt}{20pt}\selectfont\heiti\bfseries,
    number = {\chinese{section}、},
    aftername = {\hspace{1.5em}},
    beforeskip = {12pt},
    afterskip = {12pt},
  },
  subsection = {
    format = \raggedright\linespread{1.0}\fontsize{12pt}{20pt}\selectfont\heiti,
    numberformat = \rmfamily,
    aftername = {\hspace{1em}},
    beforeskip = {12pt},
    afterskip = {12pt},
  },
  subsubsection = {
    format = \raggedright\linespread{1.0}\fontsize{12pt}{20pt}\selectfont\heiti,
    numberformat = \rmfamily,
    aftername = {\hspace{1em}},
    beforeskip = {12pt},
    afterskip = {6pt},
  },
}
\fancypagestyle{fzu}{%
  \fancyhf{}
  \fancyhead{}%
  \fancyfoot[C]{\fontsize{9pt}{11pt}\selectfont 第~\thepage~页\quad 共~\pageref{LastPage}~页}%
  \renewcommand{\headrulewidth}{0pt}%
  \renewcommand{\footrulewidth}{0pt}%
}
\AtEndDocument{\label{LastPage}}
\newcommand{\proposalstep}[1]{%
  \par\textbf{#1}\par\nopagebreak[4]}
\counterwithin{figure}{section}% 图按章编号
\renewcommand{\thefigure}{\thesection-\arabic{figure}}% 形如「图 1-1」
\setcounter{tocdepth}{3}% 目录显示到三级标题

\begin{document}
\pagestyle{fzu}% 全文统一：无页眉，页脚居中页码
\linespread{1.389}\selectfont% 12pt 字号 × 1.2 × 1.389 ≈ 20 磅固定行距（FZU Body: line=400 exact）
% %=============================================================================%
%  福大「表一」封面（福州大学专业学位研究生论文开题报告）
%  版式照福大开题报告 Word 模板：
%    A4、六列竖线表格、标签宋体四号、内容宋体小四、意见栏首行缩进两字
%  列宽合计 15.60cm + 6×2\tabcolsep(3pt) + 竖线 ≈ 17.0cm（正好版心）
%
%  用法：在 proposal.tex 里把 \input{contents/cover_table} 那一行的行首 % 去掉，
%        然后把下面所有 ××× 换成自己的信息（这里的文字是直接排版的，不读变量）。
%  提示：意见栏建议打印后手写签字；如需电子版，把意见写在 \hspace{2em} 后面即可。
%=============================================================================%
\begin{titlepage}
\thispagestyle{empty}%
\begin{flushright}
{\linespread{1.0}\zihao{5}表\quad 一}\par
\end{flushright}
\vspace{-1.5em}
\begin{center}
{\linespread{1.0}\zihao{2}\heiti 福州大学专业学位研究生论文开题报告}\par
\end{center}
\vspace{-0.2em}

\begingroup
\setlength{\tabcolsep}{3pt}%
\renewcommand{\arraystretch}{1.0}%
\zihao{5}%
\noindent
\begin{tabular}{|>{\centering\arraybackslash}p{2.55cm}|>{\centering\arraybackslash}p{2.15cm}|>{\centering\arraybackslash}p{2.35cm}|>{\centering\arraybackslash}p{2.95cm}|>{\centering\arraybackslash}p{2.35cm}|>{\centering\arraybackslash}p{3.25cm}|}
\hline
{\linespread{1.0}\zihao{4}论文题目} & \multicolumn{5}{l|}{{\linespread{1.0}\zihao{-4}××××××××××××}} \\ \hline
{\linespread{1.0}\zihao{4}姓\hspace{1.1em}名} & {\linespread{1.0}\zihao{-4}×××} & {\linespread{1.0}\zihao{4}学\hspace{1.1em}号} & {\linespread{1.0}\zihao{-4}××××××××××} & {\linespread{1.0}\zihao{4}性\hspace{1.1em}别} & {\linespread{1.0}\zihao{-4}××} \\ \hline
{\linespread{1.0}\zihao{4}导\hspace{1.1em}师} & {\linespread{1.0}\zihao{-4}×××} & {\linespread{1.0}\zihao{4}学科专业} & {\linespread{1.0}\zihao{-4}××××××××} & {\linespread{1.0}\zihao{4}研究方向} & {\linespread{1.0}\zihao{-4}××××××××} \\ \hline
{\linespread{1.0}\zihao{4}学\hspace{1.1em}院} & \multicolumn{2}{c|}{{\linespread{1.0}\zihao{-4}×××××学院}} & {\linespread{1.0}\zihao{4}开题报告会时间、地点} & \multicolumn{2}{p{5.55cm}|}{\centering\zihao{-4}××××.××.××\quad ××××××××××××} \\ \hline
\parbox{\linewidth}{\linespread{1.0}\centering\zihao{4}导师审核} & \multicolumn{5}{p{13.9cm}|}{\linespread{1.0}\fontsize{11pt}{14.6pt}\selectfont%
\vspace{0.3em}%
\par\vspace{4.2em}% 留空给手写意见
\hspace{9em}导师签名：\par
\vspace{0.6em}\hspace{11em}年\qquad 月\qquad 日\par\vspace{0.3em}} \\ \hline
\multirow{2}{*}{\parbox{\linewidth}{\linespread{1.0}\centering\zihao{4}审核小组\\意见与评定成绩}} & \multicolumn{5}{p{13.9cm}|}{\linespread{1.0}\fontsize{11pt}{14.6pt}\selectfont%
\vspace{0.3em}%
\par\vspace{5.6em}% 留空给手写意见
\hspace{8em}审核小组成员签名：\par
\vspace{0.6em}\hspace{11em}年\qquad 月\qquad 日\par\vspace{0.3em}} \\ \cline{2-6}
 & \parbox{\linewidth}{\linespread{1.0}\centering\zihao{5}评定成绩\\（在相应等级打$\surd$）} & {\linespread{1.0}\zihao{5}优秀（\hspace{0.8em}）} & {\linespread{1.0}\zihao{5}良好（\hspace{0.8em}）} & {\linespread{1.0}\zihao{5}一般（\hspace{0.8em}）} & {\linespread{1.0}\zihao{5}较差（\hspace{0.8em}）} \\ \hline
\parbox{\linewidth}{\linespread{1.0}\centering\zihao{4}学位点\\意见} & \multicolumn{5}{p{13.9cm}|}{\linespread{1.0}\fontsize{11pt}{14.6pt}\selectfont%
\vspace{2.2em}%
\hspace{9em}学位点负责人签名：\par
\vspace{0.6em}\hspace{11em}年\qquad 月\qquad 日\par\vspace{0.3em}} \\ \hline
\end{tabular}
\endgroup
\end{titlepage}
% 福大「表一」封面：默认关闭；需要交表一时去掉本行行首的 %
%---------------------------------------------------------------------------%
%->> Titlepage information（封面信息）
%-> 把 ××× 换成自己的信息；表格封面关闭时，这里只作为变量保存，不参与排版
%---------------------------------------------------------------------------%
\schoollogo{scale=0.112}{images/ucas_logo}
\title{××××××××××××}% 论文中文题目
\author{×××}% 姓名
\authorid{××××××××××}% 学号
\advisor{×××}% 导师姓名
\advisortitle{×××}% 导师职称，如：研究员 / 教授
\degree{硕士}% 学位：学士、硕士、博士
\degreetype{工学}% 学位类别：理学、工学、工程、医学等
\major{×××}% 二级学科专业名称
\field{×××}% 研究方向
\institute{×××}% 培养单位
\school{福州大学}

% 说明：如果启用了福大「表一」表格封面（contents/cover_table.tex），
%       封面上的题目、姓名、学号、导师、专业等也要同步改成自己的；
%       那份表格里是直接排版的文字，不读这里的变量。
% 封面信息（表格封面关闭时，这里只提供变量，不排版）
\pagenumbering{roman}
\pagestyle{empty}% 目录页按参考稿不显示页码
%---------------------------------------------------------------------------%
%->> Frontmatter
%---------------------------------------------------------------------------%
%-
%-> 生成封面
%-
% \maketitlefujian% 简易封面（含校徽）：学校要求用「表一」表格封面，故默认关闭
%-
%-> 目录
%-
{\renewcommand{\contentsname}{目录}%
\begingroup
\ctexset{section/format={\centering\linespread{1.0}\fontsize{16pt}{20pt}\selectfont\heiti}}%
\tableofcontents% 生成目录
\endgroup}

%---------------------------------------------------------------------------%
% 目录
\clearpage
\pagenumbering{arabic}
\pagestyle{fzu}
%=============================================================================%
%  正文：福州大学开题报告五段式（一本 / 二研究内容与目标 / 三方案与可行性 /
%        四特色与创新 / 五参考文献）
%
%  写法约定
%    · 五个一级标题是学校表格里固定的，已用 \section{} 写好，不要增删、不要手打「一、」；
%      编号由模板自动生成（\chinese{section}、）。
%    · 以 % 开头的行是"填写要点"，编译时不显示；写完后留着当备忘、或者删掉都行。
%    · 需要三级标题时用 \subsubsection{}，图、表、公式、引用的写法见文末「写法示例」。
%=============================================================================%

\section{论文选题依据}

\subsection{选题背景与意义}
% 【写什么】3~4 段，只回答"为什么非做不可"，建议按这个顺序展开：
%   ① 领域/产业背景：这个测量或重建需求从哪来，有什么国家或行业需求（可先引 1~2 篇综述）；
%   ② 现有方法在本课题场景下失效的具体表现：要写现象和证据，不要只写"效果不佳"；
%   ③ 本课题的切入点，以及做成之后的价值（学术价值一句、工程价值一句）。
% 【注意】这一节不要写成文献综述，综述放到 1.2，否则两节会重复。

\subsection{国内外本学科领域的发展现状与趋势}
% 【写什么】按"技术方向"拆成 2~4 个三级标题，每个方向独立成段，段内顺序固定：
%   国外代表性方法（时间线 + 关键做法 + 关键指标） → 国内进展 → 这个方向还差什么。
%   最后回到本课题：上面这些"还差什么"，正是本课题要补的空白。
%
% 【要不要三级标题】要的话照下面这行改写，方向名自己起；不要的话保持注释即可。
% \subsubsection{方向一：×××的研究现状}
%
% 【引用】正文用 \cite{key} 引用，key 在 bib/ref.bib 里定义，上标编号由模板自动处理。


\section{论文的研究内容、研究目标及拟解决的关键问题}

\subsection{研究内容}
% 【写什么】分 3~4 条写，每条一句话说清"做什么"，不写方法细节（细节放第三章）。
%   常用写法：
%   \begin{enumerate}
%     \item 研究×××方法，解决×××问题；
%   \end{enumerate}

\subsection{研究目标}
% 【写什么】给可考核的目标：在什么数据集/场景上、达到什么精度或效率、相对现有基线提升多少。
%   反例："提高鲁棒性""实现高精度重建"——没法验收，评委会问"多少算高"。

\subsection{拟解决的关键问题}
% 【写什么】2~3 条，每条两句话：第一句说"难在哪"，第二句说"打算怎么破"。
%   这一节是评委判断你有没有想清楚的地方，宁可少写一条，也不要写空话。

% 【提示】若学院要求给出量化指标，可在 2.2 后面加一小节列张表：
%   指标名称 / 现有方法 / 本课题目标 / 验证方式。表格写法见文末「写法示例」。


\section{拟采取的研究方案及可行性分析}

\subsection{基本思路}
% 【写什么】一段话把技术路线串成因果链：输入是什么 → 中间表示是什么 → 输出是什么，
%   并说明"为什么这样串"（每一环为什么必须在前一环之后）。

\subsection{技术方案与技术路线}
% 【写什么】按研究内容拆成 3.2.1、3.2.2 … 每小节先"方案"（怎么做）再"路线"（步骤怎么走），
%   配 1~2 张流程图。方案里出现的关键公式要编号，后面正文要能引用它。

\subsection{可行性分析：研究条件与基础}
\subsubsection{已有研究的硬件基础}
% 【写什么】设备清单：型号 + 关键指标（分辨率、帧率、精度、量程），说明为什么够用。
\subsubsection{已有研究的实验基础}
% 【写什么】已经做完的事：标定结果、采集了多少组数据、跑通了什么流程、有什么初步结果。
%   这一节最好放真实数字和真实图，是评委会追问、也最容易加分的地方。

\subsection{研究进度安排}
% 【写什么】按四个阶段写，每阶段 2~4 条任务，动词开头（搭建/实现/验证/撰写）。
%   阶段之间要能接上，最后一段一定落到"写学位论文 + 答辩"。
本研究工作计划与进度时间节点安排如下所示：

\textbf{第一阶段（××××.××.××--××××.××.××）}
\begin{itemize}
	\item ×××
\end{itemize}

\textbf{第二阶段（××××.××.××--××××.××.××）}
\begin{itemize}
	\item ×××
\end{itemize}

\textbf{第三阶段（××××.××.××--××××.××.××）}
\begin{itemize}
	\item ×××
\end{itemize}

\textbf{第四阶段（××××.××.××--××××.××.××）}
\begin{itemize}
	\item ×××
\end{itemize}


\section{本课题的特色与创新之处}
% 【写什么】3~4 条，每条两段：先一句小标题，再一段"别人怎么做 → 本课题怎么做 → 差别在哪"。
%   创新点要能对应到前面的研究内容，不要凭空多出一条前面没做的事。
%
% 【写法（照结构，别照抄内容）】
% （1）×××
% 现有方法……；本课题……；因此……。

\proposalstep{（1）×××}
% 把上面的 \proposalstep 换成自己的小标题，逐条往下写。
% 若不想手动编号，也可以直接用 \paragraph{（1）×××} 或 \textbf{（1）×××}。

\proposalstep{（2）×××}

\proposalstep{（3）×××}


%=============================================================================%
%  写法示例（需要时把对应注释块复制到正文位置，再把内容换掉）
%
%  【插图】
%  \begin{figure}[H]
%    \centering
%    \includegraphics[width=0.95\textwidth]{images/proposal/你的图.jpg}
%    \caption{技术路线图}
%    \label{fig:route}
%  \end{figure}
%  正文引用：如图~\ref{fig:route}所示。      % 图号自动为「图 3-1」，按章编号
%
%  【表格】
%  \begin{table}[H]
%    \centering
%    \caption{实验条件对比}
%    \label{tab:cond}
%    \begin{tabular}{lccc}
%      \hline
%      项目 & 现有方法 & 本课题 & 备注 \\ \hline
%      ××× & ××× & ××× & ××× \\ \hline
%    \end{tabular}
%  \end{table}
%  正文引用：如表~\ref{tab:cond}所示。
%
%  【公式】
%  \begin{equation}
%    \label{eq:displacement}
%    \mathbf{d}(x) = \arg\min_{\mathbf{d}} \left\| \nabla I(x+\mathbf{d}) - \nabla I_0(x) \right\|_2^2
%  \end{equation}
%  正文引用：由式~\eqref{eq:displacement}可得……   % 公式连续编号，不带章号
%
%  【分条列举】itemize 用圆点，enumerate 用阿拉伯数字，两者都支持 \item
%
%  【加粗小标题】\textbf{×××} 或 \proposalstep{×××}（后者独立成段、不与下段分离）
%
%  【参考文献】条目写在 bib/ref.bib，正文 \cite{key} 引用；
%              bib/ref.bib 里现在只有注释掉的示例条目，取消注释、改成自己的即可。
%              注意：参考文献列表只显示正文引用过的条目，没引用就不会出现。
%=============================================================================%
% 正文

\clearpage
\renewcommand{\refname}{五、参考文献}% 参考文献作为第五条，标题不带自动编号
\addcontentsline{toc}{section}{五、参考文献}% 让目录里也出现这一条
\bibliography{bib/ref}% 参考文献（多文件版）
\end{document}
